\chapter{APPENDIX: REFERENCES}\label{appendix:references}
The text must include references to the sources you use. LUT University applies the Harvard referencing style, also called the author-date style with in-text referencing and a detailed reference list at the end.

The purpose of a reference is to provide sufficient information on a source used in the study, allowing the reader to consult the original source for further information. The reference enables the reader to find detailed information on the source easily in the list of references. You should refer to the original and most recent sources. If no new studies have been published on the topic in question, also older ones may be used.

Referring to a source means that you explain the contents of the source material in your own words. Direct citations, on the other hand, are placed in parentheses (“ ”). Plagiarism or using another person’s original material without appropriate referencing is not allowed.

\underline{Referencing technique}

In the Harvard system, the citation is placed in parentheses directly in the text to indicate the passage that has been cited from another source. Place the citation before the period that ends the sentence when it refers only to the sentence in question (Kaasinen et al. 2020, 173–174). If you are referring to more than that one sentence, introduce the source you are summarizing or paraphrasing at the beginning of the paragraph. Then, refer back to the source when needed to ensure your reader understands you are still using the same source.

Typically, the citation mentions the author (the last name is sufficient, unless authors of several sources have the same last name), the publication year and the page number. Please note that the author does not always have to be a person, but it may also be an organisation, for instance. If the source does not mention who the author is, the reference should include the name of the publication instead of the author. (Nykänen 2002, 77.) The author (or the title of the work) is very commonly mentioned as a part of a sentence: “According to a study conducted by Möttönen (2007, 68), a pike is a fish”.

If the source has more than one author, they are all mentioned in the reference by their last name and separated with the word ‘and’ or the symbol ‘\&’. Make later references to the same work with the first author’s last name and “et al.” If you reference several works published by the same author in the same year, add lower-case letters (a, b, c…) after the publication year to distinguish the sources. Use the same alphabetical organisation also in the list of references.

There are several referencing and citation styles. It is essential to consistently use the same style throughout the thesis. Examples and detailed instructions on referencing:

\href{https://libguides.lut.fi/citingelectronicdocuments}{\color{blue} \underline{LUT Academic Library’s instructions}} on how to cite electronic documents

\href{https://libguides.aalto.fi/citation_guide}{\color{blue} \underline{Aalto University citation guide}}

\href{https://www.librarydevelopment.group.shef.ac.uk/referencing/harvard.html}{\color{blue} \underline{Harvard referencing}}, University of Sheffield 